% Irodov, source PDF 78, printed p.69; 2.2 ends at original problem 61.
% One mol van der Waals equation: p = RT/(V-b)-a/V^2.
% For nu mol, U=nu*C_V*T-a*nu^2/V under constant C_V approximation.
\thermoenergyproblem{57}{69}
求一摩尔范德瓦耳斯气体在温度 $T$ 不变时，
由体积 $V_1$ 可逆等温膨胀至 $V_2$ 所做的功。
假定气体遵循
$(p+a/V^2)(V-b)=RT$。
\exsolution
将范德瓦耳斯方程解成压强形式：
\[
p(V,T)=\frac{RT}{V-b}-\frac a{V^2}.
\]
气体在可逆准静态膨胀中所做功为
\[
A=\int_{V_1}^{V_2}p(V,T)\,\mathrm dV
=RT\int_{V_1}^{V_2}\frac{\mathrm dV}{V-b}
-a\int_{V_1}^{V_2}\frac{\mathrm dV}{V^2}.
\]
分别积分，
\[
\boxed{A=
RT\ln\frac{V_2-b}{V_1-b}
+a\left(\frac1{V_2}-\frac1{V_1}\right)}.
\]
\excheck 要求 $V_1,V_2>b$；
$a\to0,b\to0$ 时恢复理想气体的
$A=RT\ln(V_2/V_1)$。

\thermoenergyproblem{58}{69}
一摩尔氧气在温度
$T=280\,\mathrm K$ 不变时，
由体积 $V_1=1.00\,\mathrm L$
膨胀至 $V_2=5.0\,\mathrm L$。
把氧气视为范德瓦耳斯气体，求
(a) 内能变化 $\Delta U$；
(b) 气体吸热 $Q$。
\exsolution
在常定容热容的范德瓦耳斯模型中，
一摩尔内能可写成
\[
U(T,V)=C_VT-\frac aV+\text{常数}.
\]
等温过程中 $\Delta T=0$，
所以
\[
\boxed{\Delta U=a\left(\frac1{V_1}-\frac1{V_2}\right)}.
\]
另一方面第57题给气体做功
\[
A=RT\ln\frac{V_2-b}{V_1-b}
+a\left(\frac1{V_2}-\frac1{V_1}\right).
\]
由第一定律 $Q=\Delta U+A$，
两项中含 $a$ 的部分相消：
\[
\boxed{Q=RT\ln\frac{V_2-b}{V_1-b}}.
\]
取氧气常用范德瓦耳斯参数
$a\simeq1.36\,\mathrm{atm\,L^2/mol^2}$、
$b\simeq0.0318\,\mathrm{L/mol}$：
\[
\Delta U\simeq1.09\,\mathrm{atm\,L}
\simeq\boxed{0.11\,\mathrm{kJ}},
\quad
Q\simeq(8.314)(280)
\ln\frac{5-0.0318}{1-0.0318}
\simeq\boxed{3.8\,\mathrm{kJ}}.
\]
\excheck 与理想气体等温过程的 $\Delta U=0$ 不同：
分子间吸引项 $-a/V$ 导致范德瓦耳斯内能
还依赖体积。

\thermoenergyproblem{59}{69}
范德瓦耳斯气体满足
$(p+a/V^2)(V-b)=RT$（一摩尔）。
在定容摩尔热容 $C_V$ 近似为常数的条件下，求
(a) 可逆绝热曲线在 $T,V$ 坐标中的方程；
(b) 热容差 $C_p-C_V$ 随 $T,V$ 的函数。
\exsolution
由范德瓦耳斯内能
$U=C_VT-a/V$，
有
$\mathrm dU=C_V\,\mathrm dT+(a/V^2)\mathrm dV$。
可逆绝热 $\delta Q=0$，
第一定律给
\[
0=C_V\,\mathrm dT+
\left(\frac a{V^2}+p\right)\mathrm dV
=C_V\,\mathrm dT+\frac{RT}{V-b}\mathrm dV.
\]
分离变量积分：
\[
C_V\frac{\mathrm dT}{T}
=-R\frac{\mathrm dV}{V-b}
\quad\Longrightarrow\quad
\boxed{T(V-b)^{R/C_V}=\text{常数}}.
\]
(b) 使用一般热容差恒等式
\[
C_p-C_V=
-T\frac{(\partial p/\partial T)_V^2}
{(\partial p/\partial V)_T}.
\]
由于
\[
\left(\frac{\partial p}{\partial T}\right)_V=\frac R{V-b},
\qquad
\left(\frac{\partial p}{\partial V}\right)_T
=-\frac{RT}{(V-b)^2}+\frac{2a}{V^3},
\]
得
\[
\boxed{C_p-C_V=
\frac R{1-\dfrac{2a(V-b)^2}{RTV^3}}}.
\]
\excheck 在稀薄气体极限
$a,b\to0$，两结果退化为
$TV^{R/C_V}=\text{常数}$ 和 $C_p-C_V=R$。
热容差表达式仅在稳定的单相分支有物理意义。

\thermoenergyproblem{60}{69}
两个绝热刚性容器由阀门相通：
第一容器体积 $V_1=10\,\mathrm L$，
内有 $\nu=2.5$ mol 二氧化碳；
第二容器体积 $V_2=100\,\mathrm L$，
原来是真空。打开阀门后气体自由膨胀，
假设遵从范德瓦耳斯方程，
求气体的温度变化 $\Delta T$。
\exsolution
两容器总系统绝热且刚性，
向真空自由膨胀不对外做功，
因此总内能守恒：
\[
\Delta U=0.
\]
对 $\nu$ 摩尔的范德瓦耳斯模型，
内能为
\[
U=\nu C_VT-\frac{a\nu^2}{V}+\text{常数}.
\]
最终容积是\emph{两容器之和}
$V_f=V_1+V_2$。
故
\[
0=\nu C_V(T_f-T_i)
-a\nu^2\left(\frac1{V_f}-\frac1{V_1}\right).
\]
因此
\[
\boxed{\Delta T=T_f-T_i=
\frac{a\nu}{C_V}
\left(\frac1{V_1+V_2}-\frac1{V_1}\right)<0}.
\]
采用二氧化碳参数
$a\simeq3.59\,\mathrm{atm\,L^2/mol^2}$，
并与原书此温区热容假设一致地取
$\gamma\simeq1.30$、$C_V=R/(\gamma-1)$，
则
\[
\boxed{\Delta T\simeq-3.0\,\mathrm K}.
\]
\excheck 理想气体向真空自由膨胀温度不变，
但范德瓦耳斯吸引作用导致内能中的
$-a\nu^2/V$ 随体积改变，温度因而降低。
若采用\emph{不同温区的刚性线性分子近似}
$C_V=5R/2$，则会得约 $-4\,\mathrm K$；
因此数值取决于二氧化碳的热容模型，
解析式不依赖这一近似选择。

\thermoenergyproblem{61}{69}
$\nu=3.0$ mol 二氧化碳
在\emph{向真空自由膨胀}时，
由体积 $V_1=5.0\,\mathrm L$
变为 $V_2=10\,\mathrm L$，
欲保持气体温度不变，
需要向气体传入多少热量？
按范德瓦耳斯气体处理。
\exsolution
气体向真空自由膨胀，
对外做功 $A=0$。
范德瓦耳斯总内能
$U=\nu C_VT-a\nu^2/V$，
在 $\Delta T=0$ 时：
\[
\Delta U=
a\nu^2\left(\frac1{V_1}-\frac1{V_2}\right).
\]
由第一定律
$Q=\Delta U+A=\Delta U$，得到
\[
\boxed{Q=a\nu^2\frac{V_2-V_1}{V_1V_2}}.
\]
用 $a\simeq3.59\,\mathrm{atm\,L^2/mol^2}$，
\[
Q\simeq3.59(3.0)^2
\left(\frac15-\frac1{10}\right)
\,\mathrm{atm\,L}
=3.23\,\mathrm{atm\,L}
\simeq\boxed{0.33\,\mathrm{kJ}}.
\]
\excheck \emph{向真空}的自由膨胀与
第57题\emph{可逆等温}膨胀不同。
在自由膨胀中没有准静态 $p\,\mathrm dV$
外功，不能套用第57题的做功积分。
